% Attention !! Rien dans le chapitre ici sur les projections parallèles : à rajouter ?!
%
%

\chapter{Proportionnalité et projections parallèles} 

\section{Tableaux de proportionnalité}

Un petit vendeur de pâtisseries vend ses croissants à $1,5$ EURO pièce.

\begin{tasks}
  \task Combien vas-tu payer pour 6 croissants?
  \task Complète le tableau suivant.
\end{tasks} 
  \begin{table}[]
    \resizebox{\textwidth}{!}{%
    \begin{tabular}{llllll}
    Nombre de croissants & 0 & 1 &   & 18 &    \\
    Prix (en euro)          &   &   & 6 &    & 54
    \end{tabular}%
    }
  \end{table}
%De nouveau le tableau est mal mis sur la page je comprends pas pourquoi

Deux grandeurs proportionnelles sont \dots

\begin{exercise}
  Vrai ou faux? 
  \begin{tasks}
    \task Le nombre de bouteilles d'eau achetées et le prix payé sont proportionnels \dots
    \task La taille d'une personne est proportionnelle à son âge \dots
    \task Le salaire d'un professeur et le nombre de ses élèves en classe sont proportionnels \dots
    \task Le prix d'un livre et le nombre de pages qu'il contient sont proportionnels \dots
    \task Le numéro d'un épisode de ta série préférée est proportionnel au temps de cet épisode \dots
    \task Le périmètre d'un carré est proportionnel à la mesure de son côté \dots
    \task Le temps de trajet sur l'autoroute est proportionnel à la vitesse du véhicule \dots
  \end{tasks}
\end{exercise}

\begin{exercise}
En t'aidant de la définition des grandeurs proportionnelles, complète les tableaux ci-dessous en indiquant sur les flèches les opérations effectuées.
%Refaire les tableaux avec flèches !! 
\end{exercise}

Le coefficient de proportionnalité $k$ entre deux grandeurs
proportionnelles $x$ et $y$ est \dots

\exemple 
%tableau à faire 

\begin{exercise}
  Détermine le coefficient de proportionnalité reliant les 2 grandeurs de chaue tableau de proportionnalité ci-dessous, et utilise-le pour compléter les cases vides.
%tableaux à refaire
\end{exercise}

\begin{exercise}
  Dans chaque cas, détermine si les grandeurs $x$ et $y$ sont proportionnelles. Si oui, détermine le coefficient de proportionnalité.
  %Tableaux à refaire
\end{exercise}

\begin{exercise}
  Complète les tableaux suivants grâce au coefficient de proportionnalité.
  %Tableaux à refaire
\end{exercise}

\section{Règle de 3 et proportionnalité }

\begin{exercise}
  Marie a acheté 4 Mangas de OnePiece et elle a payé 20euros. \\
  Combien aurait-elle payé pour 6 albums? Ecris tous tes calculs.
\end{exercise}

\begin{exercise}
  Ma voiture consomme, en moyenne, 4 litress de carburant pour 100 km. Quelle distance pourrais-je parcourir avec un plein de 50 litres? 
\end{exercise}

\begin{exercise}
  Chez Colruyt, un lot de 4 bouteilles de vin est affiché à 22euros. On peut acheter les bouteilles à la pièce. \\
  Combien vais-je payer pour l'achat de 10 bouteilles? 
\end{exercise}

\begin{exercise}
  Pour une exposition, un groupe de 12 élèves a payé 84 euros. Qu'auraient payé 17 élèves?
\end{exercise}

\begin{exercise}
  Sur un plan au $\dfrac{1}{2000}$, un champ rectangulaire ests représenté par un rectangle mesurant 2cm de long et 1,4cm de large. \\
  Quelles sont les dimensions réelles (en mètres)  de ce champ ? \\
  Quelle est l'aire de ce champ (en mètres carrés)?
\end{exercise}

\section{Tableaux, graphiques}

\begin{exercise}
  Voici les tarifs d'entrée à la piscine de la commune. 
  %Tableau à faire
  \begin{tasks}
    \task Est-ce un tableau de proportionnalité? Justifie.
    \task Construis le graphique évolutif qui traduit la situation. 
    %Mettre un quadrillage et des axes pour le graphique
  \end{tasks}
\end{exercise}

Note : Le graphique qui traduit une situation de proportionnalité est \textbf{toujours} une droite passant par l'origine du repère.

\begin{exercise}
  Parmi les graphiques suivants, quel est celui qui traduit une situation de proportionnalité? Justifie en disant pourquoi les \textbf{autres} ne le sont pas.
  %ajouter graphiques
\end{exercise}

\section{Exercices mélangés}
\begin{exercise}
  Exercice issu d'un CE1D.
  %Image à remettre
\end{exercise}

\begin{exercise}
  Pour télécharger 3 chansons sur internet, à une certaine époque, il fallait 1 minute.
  \begin{tasks}
    \task Complète, en te basant sur ce temps moyen de téléchargement, le tableau de proportionnalité suivant : \\
    \begin{tabularx}{\linewidth}{XX}
      \toprule
      \thead{Nombre de chansons}         & \thead{Durée de téléchargement (en secondes)}  \\
      \midrule
        &  120\\
      9 &     \\
        & 500   \\
      \bottomrule
    \end{tabularx}

    \task Calcule le nombre de chansons que tu pourrais télécharger, à la même vitesse, en une demi-heure.
  \end{tasks}
\end{exercise}

\begin{exercise}
  Un robinet ouvert à fond remplit un seau de 12 litres en 18 secondes. Etudie le temps qu'il faudra pour remplir des seaux de différentes capacités.
  \begin{tasks}
    \task Complète le tableau suivant correspondant à la situation donnée : \\
    % Faire le tableau
    \task Trace le graphique liant la quantité d'eau déversée par le robient en fonction du temps.
    %Mettre un quadrillage
    \task Justifie à l'aide du tableau \textbf{et} à l'aide du graphique qu'il s'agit d'une situation de proportionnalité.
  \end{tasks}
\end{exercise}

\begin{exercise}
  La hauteur $[AC]$ du triangle $BDC$ mesure 2cm. Le segment $[AB]$ mesure 3cm. \\
  Construis un agrandissement de la figure si tu sais que le segment $[A'B']$ devra mesurer, après l'agrandissement, 4,5cm.
  %Insérer l'image du triangle.
\end{exercise}

\begin{exercise}
  Construis le polygone $A'B'C'D'$, semblable au polygone $ABCD$ tel que $k=\dfrac{1}{2}$.
  %Insérer image du polygone
\end{exercise}

\begin{exercise}
  Dans une école de  Wallonie,  en  1ère  année,  il  y  a 120 élèves.  Il en  a 30 qui sont en immersion, 25 qui suivent le latin, 40 qui font l’activité maths et le reste qui sont des doubleurs. Représente la répartition des élèves sur un diagramme circulaire.
  \begin{table}[]
    \resizebox{\textwidth}{!}{%
    \begin{tabular}{|l|l|l|l|}
    \hline
    \textbf{Options} & \textbf{Nombre d'élèves} & \textbf{Pourcentage} & \textbf{Amplitude} \\\hline
    Immersion        &                          &                      &                    \\ \hline
    Latin            &                          &                      &                    \\ \hline
    Activité de math &                          &                      &                    \\ \hline
    Doubleurs        &                          &                      &                    \\ \hline
    \textbf{Total}   &                          &                      &    \\ \hline               
    \end{tabular}%
    }
    \end{table}
    % de nouveau le tableau ne se met pas là où je veux qu'il se mette... pas cool ! 
    %Rajouter un cercle avec centre pour leur diagramme 
\end{exercise}